% ============================================================
\chapter{Fonctions Mesurables}
\label{ch:fonctions_mesurables}
% ============================================================

Avant de pouvoir int\'egrer, il faut savoir quelles fonctions ont le droit d'\^etre int\'egr\'ees. C'est la question de la \emph{mesurabilit\'e}~: une fonction est mesurable si les images r\'eciproques des ensembles \og raisonnables \fg sont elles-m\^emes mesurables. Cette condition, qui semble technique, est en r\'ealit\'e tr\`es naturelle~: elle garantit que l'on peut d\'efinir sans ambigu\"it\'e la probabilit\'e que $f(X)$ tombe dans un intervalle donn\'e, ou l'int\'egrale de $f$ par rapport \`a une mesure. Les fonctions continues sont toujours mesurables, mais la classe des fonctions mesurables est consid\'erablement plus vaste --- c'est pr\'ecis\'ement cette souplesse qui rend l'int\'egrale de Lebesgue si puissante.

\section{D\'efinitions et crit\`eres de mesurabilit\'e}

\begin{definition}[Fonction mesurable]
Soit $(X, \mathcal{A})$ un espace mesurable. Une fonction
$f : X \to \overline{\R} = [-\infty, +\infty]$ est dite
\emph{$\mathcal{A}$-mesurable} (ou simplement \emph{mesurable}) si~:
\[
\forall\, a \in \R, \quad f^{-1}([-\infty, a)) = \{x \in X : f(x) < a\} \in
\mathcal{A}.
\]
\end{definition}

\begin{proposition}[Crit\`eres \'equivalents]
\label{prop:criteres_mesurabilite}
Pour $f : X \to \overline{\R}$, les conditions suivantes sont \'equivalentes~:
\begin{enumerate}
    \item $\{f < a\} \in \mathcal{A}$ pour tout $a \in \R$\,;
    \item $\{f \leq a\} \in \mathcal{A}$ pour tout $a \in \R$\,;
    \item $\{f > a\} \in \mathcal{A}$ pour tout $a \in \R$\,;
    \item $\{f \geq a\} \in \mathcal{A}$ pour tout $a \in \R$.
\end{enumerate}
De plus, si l'une de ces conditions est v\'erifi\'ee, alors
$\{f = a\} \in \mathcal{A}$, $\{f = +\infty\} \in \mathcal{A}$ et
$\{f = -\infty\} \in \mathcal{A}$.
\end{proposition}

\begin{proof}
$(1) \Rightarrow (2)$~: $\{f \leq a\} = \bigcap_{n=1}^{\infty} \{f < a + 1/n\}$.

$(2) \Rightarrow (3)$~: $\{f > a\} = \{f \leq a\}^c$.

$(3) \Rightarrow (4)$~: $\{f \geq a\} = \bigcap_{n=1}^{\infty} \{f > a - 1/n\}$.

$(4) \Rightarrow (1)$~: $\{f < a\} = \{f \geq a\}^c$.

Enfin, $\{f = a\} = \{f \leq a\} \cap \{f \geq a\}$,
$\{f = +\infty\} = \bigcap_n \{f > n\}$, $\{f = -\infty\} = \bigcap_n \{f < -n\}$.
\end{proof}

\begin{exemple}
\begin{enumerate}
    \item Toute fonction constante est mesurable.
    \item La fonction indicatrice $\mathbf{1}_A$ est mesurable si et seulement si
    $A \in \mathcal{A}$.
    \item Si $X$ est un espace topologique et $\mathcal{A} = \mathcal{B}(X)$, toute
    fonction continue $f : X \to \R$ est mesurable.
    \item Toute fonction monotone $f : \R \to \R$ est bor\'elienne.
\end{enumerate}
\end{exemple}

\section{Op\'erations sur les fonctions mesurables}

\begin{theoreme}[Stabilit\'e alg\'ebrique]
\label{thm:stabilite_algebrique}
Soient $f, g : X \to \R$ deux fonctions mesurables et $\alpha \in \R$. Alors~:
\begin{enumerate}
    \item $\alpha f$ est mesurable\,;
    \item $f + g$ est mesurable (si bien d\'efinie)\,;
    \item $fg$ est mesurable\,;
    \item $\abs{f}$ est mesurable\,;
    \item $f/g$ est mesurable (l\`a o\`u $g \neq 0$)\,;
    \item $f^+$, $f^-$, $\max(f, g)$, $\min(f, g)$ sont mesurables.
\end{enumerate}
\end{theoreme}

\begin{proof}
(2) L'id\'ee cl\'e est l'observation~:
\[
\{f + g < a\} = \bigcup_{r \in \Q} (\{f < r\} \cap \{g < a - r\}).
\]
C'est une r\'eunion d\'enombrable d'ensembles mesurables, donc mesurable.

(3) On utilise l'identit\'e $fg = \frac{1}{4}[(f+g)^2 - (f-g)^2]$ et le fait que
$h^2$ est mesurable si $h$ l'est (car $\{h^2 < a\} = \{-\sqrt{a} < h < \sqrt{a}\}$
pour $a > 0$).

(4) $\{\abs{f} < a\} = \{-a < f < a\} = \{f > -a\} \cap \{f < a\}$.

(6) $\max(f, g) = \frac{1}{2}(f + g + \abs{f - g})$ et
$\min(f, g) = \frac{1}{2}(f + g - \abs{f - g})$.
\end{proof}

\section{Limites de fonctions mesurables}

\begin{theoreme}[Stabilit\'e par passage \`a la limite]
\label{thm:stabilite_limite}
Soit $(f_n)_{n \in \N}$ une suite de fonctions mesurables $X \to \overline{\R}$.
Alors les fonctions suivantes sont mesurables~:
\begin{enumerate}
    \item $\sup_n f_n$ et $\inf_n f_n$\,;
    \item $\limsup_n f_n$ et $\liminf_n f_n$\,;
    \item si $f = \lim_n f_n$ existe (ponctuellement), alors $f$ est mesurable.
\end{enumerate}
\end{theoreme}

\begin{proof}
(1) $\{\sup_n f_n > a\} = \bigcup_n \{f_n > a\} \in \mathcal{A}$.
De m\^eme, $\{\inf_n f_n < a\} = \bigcup_n \{f_n < a\} \in \mathcal{A}$.

(2) $\limsup_n f_n = \inf_n \sup_{k \geq n} f_k$ et
$\liminf_n f_n = \sup_n \inf_{k \geq n} f_k$, donc mesurables par (1).

(3) Si $f = \lim_n f_n$, alors $f = \limsup_n f_n = \liminf_n f_n$.
\end{proof}

\begin{attention}[title=\textbf{Contraste avec la continuit\'e}]
Une limite ponctuelle de fonctions continues n'est pas n\'ecessairement continue
(par exemple, $f_n(x) = x^n$ sur $[0,1]$). Mais une limite ponctuelle de fonctions
mesurables est \emph{toujours} mesurable. C'est un avantage majeur de la
mesurabilit\'e sur la continuit\'e.
\end{attention}

\section{Fonctions simples (\'etag\'ees)}

\begin{definition}[Fonction simple]
Une fonction $\varphi : X \to \R$ est dite \emph{simple} (ou \emph{\'etag\'ee}) si
elle ne prend qu'un nombre fini de valeurs. Elle s'\'ecrit~:
\[
\varphi = \sum_{k=1}^m c_k \mathbf{1}_{A_k},
\]
o\`u $c_1, \ldots, c_m$ sont les valeurs distinctes prises par $\varphi$ et
$A_k = \varphi^{-1}(\{c_k\})$ sont les pr\'eimages correspondantes.
$\varphi$ est mesurable si et seulement si chaque $A_k \in \mathcal{A}$.
\end{definition}

\begin{theoreme}[Approximation par fonctions simples]
\label{thm:approximation_simples}
Soit $f : X \to [0, +\infty]$ une fonction mesurable. Il existe une suite
$(\varphi_n)_{n \geq 1}$ de fonctions simples mesurables telle que~:
\begin{enumerate}
    \item $0 \leq \varphi_1 \leq \varphi_2 \leq \ldots$\,;
    \item $\varphi_n(x) \uparrow f(x)$ pour tout $x \in X$\,;
    \item si $f$ est born\'ee, la convergence est uniforme.
\end{enumerate}
Plus g\'en\'eralement, toute fonction mesurable $f : X \to \overline{\R}$ est limite
ponctuelle d'une suite de fonctions simples mesurables.
\end{theoreme}

\begin{proof}
Pour $n \geq 1$ et $x \in X$, posons~:
\[
\varphi_n(x) = \begin{cases}
k \cdot 2^{-n} & \text{si } k \cdot 2^{-n} \leq f(x) < (k+1) \cdot 2^{-n},
\quad k = 0, 1, \ldots, n \cdot 2^n - 1, \\
n & \text{si } f(x) \geq n.
\end{cases}
\]
Cette fonction est simple (prenant au plus $n \cdot 2^n + 1$ valeurs) et mesurable
(chaque pr\'eimage est un ensemble mesurable de la forme $\{k \cdot 2^{-n} \leq
f < (k+1) \cdot 2^{-n}\}$).

\emph{Monotonie}~: pour tout $x$, $\varphi_n(x) \leq \varphi_{n+1}(x)$ car le
d\'ecoupage au rang $n+1$ est un raffinement de celui au rang $n$.

\emph{Convergence}~: si $f(x) < \infty$, pour $n > f(x)$, on a
$0 \leq f(x) - \varphi_n(x) < 2^{-n} \to 0$. Si $f(x) = +\infty$,
$\varphi_n(x) = n \to +\infty$.

\emph{Convergence uniforme si $f$ born\'ee}~: si $f \leq M$, pour $n > M$,
$\sup_x \abs{f(x) - \varphi_n(x)} \leq 2^{-n}$.

Pour $f$ g\'en\'erale, on \'ecrit $f = f^+ - f^-$ et on approche $f^+$ et $f^-$
s\'epar\'ement.
\end{proof}

\begin{formules}[title=\textbf{Construction de l'approximation}]
\[
\varphi_n = \sum_{k=0}^{n2^n - 1} \frac{k}{2^n}\, \mathbf{1}_{\{k/2^n \leq f <
(k+1)/2^n\}} + n\, \mathbf{1}_{\{f \geq n\}}.
\]
\end{formules}

\section{Fonctions mesurables et fonctions bor\'eliennes}

\begin{proposition}
Soit $(X, \mathcal{A})$ un espace mesurable. Si $f : X \to \R$ est mesurable et
$g : \R \to \R$ est bor\'elienne, alors $g \circ f$ est mesurable.
\end{proposition}

\begin{proof}
Pour tout $B \in \mathcal{B}(\R)$, $(g \circ f)^{-1}(B) = f^{-1}(g^{-1}(B))$.
Comme $g$ est bor\'elienne, $g^{-1}(B) \in \mathcal{B}(\R)$. Comme $f$ est
$(\mathcal{A}, \mathcal{B}(\R))$-mesurable, $f^{-1}(g^{-1}(B)) \in \mathcal{A}$.
\end{proof}

\begin{corollaire}
Si $f$ est mesurable, alors $e^f$, $\sin(f)$, $f^p$ ($p > 0$), $\log(\abs{f})$
(l\`a o\`u $f \neq 0$) sont mesurables.
\end{corollaire}

\section{Convergence presque partout et en mesure}

\begin{definition}[Convergence presque partout]
On dit que $f_n \to f$ \emph{$\mu$-presque partout} ($\mu$-p.p.) si
\[
\mu\!\left(\{x \in X : f_n(x) \not\to f(x)\}\right) = 0.
\]
\end{definition}

\begin{definition}[Convergence en mesure]
On dit que $f_n \to f$ \emph{en mesure} si pour tout $\varepsilon > 0$~:
\[
\mu\!\left(\{x \in X : \abs{f_n(x) - f(x)} > \varepsilon\}\right) \to 0
\quad (n \to \infty).
\]
\end{definition}

\begin{theoreme}
Si $\mu(X) < \infty$ et $f_n \to f$ p.p., alors $f_n \to f$ en mesure.
\end{theoreme}

\begin{proof}
Posons $A_n(\varepsilon) = \{x : \abs{f_n(x) - f(x)} > \varepsilon\}$ et
$B_N(\varepsilon) = \bigcup_{n \geq N} A_n(\varepsilon)$. Alors
$B_N(\varepsilon) \downarrow \{x : \abs{f_n(x) - f(x)} > \varepsilon \text{ pour
une infinit\'e de } n\} \subset \{f_n \not\to f\}$.

Par hypoth\`ese, $\mu(\{f_n \not\to f\}) = 0$, donc $\mu(B_N(\varepsilon)) \to 0$.
Comme $A_N(\varepsilon) \subset B_N(\varepsilon)$,
$\mu(A_N(\varepsilon)) \to 0$.
\end{proof}

\begin{theoreme}[R\'eciproque partielle]
Si $f_n \to f$ en mesure, il existe une sous-suite $(f_{n_k})$ telle que
$f_{n_k} \to f$ p.p.
\end{theoreme}

\begin{proof}
Pour chaque $k$, choisir $n_k$ tel que
$\mu(\{\abs{f_{n_k} - f} > 2^{-k}\}) < 2^{-k}$. Par le lemme de Borel--Cantelli
(lemme~\ref{lem:borel_cantelli}), $\mu(\limsup_k \{\abs{f_{n_k} - f} > 2^{-k}\})
= 0$. Hors de cet ensemble, $\abs{f_{n_k} - f} \leq 2^{-k}$ pour $k$ assez grand,
donc $f_{n_k} \to f$.
\end{proof}

\begin{theoreme}[Egorov]
\label{thm:egorov}
Soit $\mu(X) < \infty$ et $f_n \to f$ $\mu$-p.p. Pour tout $\delta > 0$, il existe
$A \in \mathcal{A}$ avec $\mu(A^c) < \delta$ et $f_n \to f$ uniform\'ement sur $A$.
\end{theoreme}

\begin{proof}
Pour $m, n \in \N^*$, posons $E_{m,n} = \bigcup_{k \geq n} \{|f_k - f| \geq 1/m\}$.
Pour $m$ fix\'e, $E_{m,n} \downarrow E_m$ o\`u $E_m \subset \{f_n \not\to f\}$,
donc $\mu(E_m) = 0$. Ainsi $\mu(E_{m,n}) \to 0$ quand $n \to \infty$.

Choisissons $n_m$ tel que $\mu(E_{m,n_m}) < \delta/2^m$. Posons
$A = X \setminus \bigcup_m E_{m,n_m}$. Alors
$\mu(A^c) \leq \sum_m \delta/2^m = \delta$. Sur $A$, pour tout $m$, et tout
$k \geq n_m$, $|f_k - f| < 1/m$, ce qui est la convergence uniforme.
\end{proof}

\section{Th\'eor\`eme de Lusin}

\begin{theoreme}[Lusin]
\label{thm:lusin}
Soit $f : \R^n \to \R$ une fonction Lebesgue-mesurable et $\varepsilon > 0$.
Il existe un ferm\'e $F \subset \R^n$ tel que $\lambda^n(F^c) < \varepsilon$ et
$f|_F$ est continue.
\end{theoreme}

\begin{intuition}[title=\textbf{Interpr\'etation de Lusin}]
\og Toute fonction mesurable est presque continue. \fg Plus pr\'ecis\'ement,
on peut rendre une fonction mesurable continue en retirant un ensemble de mesure
arbitrairement petite.
\end{intuition}

\section{Exercices}

\begin{exercice}
Montrer que si $f : \R \to \R$ est monotone, alors $f$ est bor\'elienne.
\end{exercice}

\begin{exercice}
Soit $f : X \to \R$ mesurable et $g : X \to \R$ telle que $f = g$ $\mu$-p.p.
Montrer que $g$ n'est pas n\'ecessairement mesurable (donner un contre-exemple).
En revanche, montrer que si $\mu$ est compl\`ete, alors $g$ est mesurable.
\end{exercice}

\begin{exercice}
Montrer que la convergence en mesure est m\'etrisable lorsque $\mu(X) < \infty$
par la distance $d(f, g) = \int_X \frac{\abs{f - g}}{1 + \abs{f - g}}\, d\mu$.
\end{exercice}

\begin{exercice}
Construire une suite $(f_n)$ sur $([0,1], \lambda)$ qui converge en mesure vers $0$
mais ne converge en aucun point. (\emph{Fonctions typographes.})
\end{exercice}

\begin{exercice}
Soit $(f_n)$ une suite de fonctions mesurables positives. Montrer que
$\{x : \sum_n f_n(x) < \infty\}$ est un ensemble mesurable.
\end{exercice}
